% GENERATED FILE. Edit canonical lessons, then run npm run generate:books.
% canonical-source-hash: fnv1a64:9d16d699c6409be7
% canonical-lessons: HI-C139-matar, HI-C139-khira, HI-C139-am, HI-C139-kela, HI-C139-seb

\chapter{Peas, Cucumber, Mango, Banana, Apple}
\label{ch:hi-peas-cucumber-mango-banana-apple}
\hlchaptermodality{139}

\begin{chapteropening}
\textbf{By the end of this chapter:} \emph{I can say peas, a cucumber, a mango, a banana and an apple.}

Complete the last lesson of chapter 139: i can say peas, a cucumber, a mango, a banana and an apple.
\end{chapteropening}

\section[maṭar]{\hi{मटर} (maṭar) — peas}

\label{lesson:HI-C139-matar}



\begin{quote}
Before the new one: say the Hindi for a carrot, then the Hindi for a cauliflower.
\end{quote}

\subsection*{You'll want to know: \hi{मटर}}
\textbf{\hi{मटर}} — \emph{maṭar} — \textquotedblleft{}peas\textquotedblright{}.

\begin{grammarlens}[title={What you've built}]
One more word for this chapter.
\end{grammarlens}

\subsection*{Your turn}
\begin{itemize}
\raggedright
  \item \emph{Say it:} \emph{maṭar}
  \item \emph{Say it:} \emph{maṭar}, once more
  \item \emph{Say it:} say \emph{maṭar}
  \item \emph{Recall:} say the Hindi for a carrot, then the Hindi for a cauliflower, then say \emph{maṭar} again
\end{itemize}

\begin{tcolorbox}[breakable,colback=teal!4,colframe=teal!35!black,title={Before you move on}]
What does \hi{मटर} mean? (\textquotedblleft{}Peas\textquotedblright{}.) Say it once more.
\end{tcolorbox}
\section[khīrā]{\hi{खीरा} (khīrā) — a cucumber}

\label{lesson:HI-C139-khira}



\begin{quote}
Before the new one: say the Hindi for a cauliflower, then the Hindi for peas.
\end{quote}

\subsection*{You'll want to know: \hi{खीरा}}
\textbf{\hi{खीरा}} — \emph{khīrā} — \textquotedblleft{}a cucumber\textquotedblright{}.

\begin{grammarlens}[title={What you've built}]
One more word for this chapter.
\end{grammarlens}

\subsection*{Your turn}
\begin{itemize}
\raggedright
  \item \emph{Say it:} \emph{khīrā}
  \item \emph{Say it:} \emph{khīrā}, once more
  \item \emph{Say it:} say \emph{khīrā}
  \item \emph{Recall:} say the Hindi for a cauliflower, then the Hindi for peas, then say \emph{khīrā} again
\end{itemize}

\begin{tcolorbox}[breakable,colback=teal!4,colframe=teal!35!black,title={Before you move on}]
What does \hi{खीरा} mean? (\textquotedblleft{}A cucumber\textquotedblright{}.) Say it once more.
\end{tcolorbox}
\section[ām]{\hi{आम} (ām) — a mango}

\label{lesson:HI-C139-am}



\begin{quote}
Before the new one: say the Hindi for peas, then the Hindi for a cucumber.

Say \textbf{\hi{मुझे} \hi{चाय} \hi{पसंद} \hi{है}}. \textbf{\hi{पसंद}} is not a verb, and \textbf{\hi{मुझे}} means \emph{to me}.
\end{quote}

\subsection*{You'll want to know: \hi{आम}}
\textbf{\hi{आम}} — \emph{ām} — \textquotedblleft{}a mango\textquotedblright{}.

\begin{grammarlens}[title={What you've built}]
One more word for this chapter.
\end{grammarlens}

\subsection*{Your turn}
\begin{itemize}
\raggedright
  \item \emph{Say it:} \emph{ām}
  \item \emph{Say it:} \emph{ām}, once more
  \item \emph{Say it:} say \emph{ām}
  \item \emph{Recall:} say the Hindi for peas, then the Hindi for a cucumber, then say \emph{ām} again
\end{itemize}

\begin{tcolorbox}[breakable,colback=teal!4,colframe=teal!35!black,title={Before you move on}]
What does \hi{आम} mean? (\textquotedblleft{}A mango\textquotedblright{}.) Say it once more.
\end{tcolorbox}
\section[kelā]{\hi{केला} (kelā) — a banana}

\label{lesson:HI-C139-kela}



\begin{quote}
Before the new one: say the Hindi for a cucumber, then the Hindi for a mango.
\end{quote}

\subsection*{You'll want to know: \hi{केला}}
\textbf{\hi{केला}} — \emph{kelā} — \textquotedblleft{}a banana\textquotedblright{}.

\begin{grammarlens}[title={What you've built}]
One more word for this chapter.
\end{grammarlens}

\subsection*{Your turn}
\begin{itemize}
\raggedright
  \item \emph{Say it:} \emph{kelā}
  \item \emph{Say it:} \emph{kelā}, once more
  \item \emph{Say it:} say \emph{kelā}
  \item \emph{Recall:} say the Hindi for a cucumber, then the Hindi for a mango, then say \emph{kelā} again
\end{itemize}

\begin{tcolorbox}[breakable,colback=teal!4,colframe=teal!35!black,title={Before you move on}]
What does \hi{केला} mean? (\textquotedblleft{}A banana\textquotedblright{}.) Say it once more.
\end{tcolorbox}
\section[seb]{\hi{सेब} (seb) — an apple}

\label{lesson:HI-C139-seb}



\begin{quote}
Before the new one: say the Hindi for a mango, then the Hindi for a banana.
\end{quote}

\subsection*{You'll want to know: \hi{सेब}}
\textbf{\hi{सेब}} — \emph{seb} — \textquotedblleft{}an apple\textquotedblright{}.

\begin{grammarlens}[title={What you've built}]
One more word for this chapter.
\end{grammarlens}

\subsection*{Your turn}
\begin{itemize}
\raggedright
  \item \emph{Say it:} \emph{seb}
  \item \emph{Say it:} \emph{seb}, once more
  \item \emph{Say it:} say \emph{seb}
  \item \emph{Recall:} say the Hindi for a mango, then the Hindi for a banana, then say \emph{seb} again
\end{itemize}

\begin{tcolorbox}[breakable,colback=teal!4,colframe=teal!35!black,title={Before you move on}]
What does \hi{सेब} mean? (\textquotedblleft{}An apple\textquotedblright{}.) Say it once more.
\end{tcolorbox}
